\documentclass{article}
\usepackage[utf8]{inputenc}
\usepackage[T1]{fontenc}
\usepackage{geometry}
\usepackage{graphicx}
\usepackage{float}
\usepackage{hyperref}
\usepackage[english,indonesian]{babel}

\title{Optimasi Pemrosesan Template LaTeX Menggunakan Pendekatan Berbasis AI}
\author{Academic Support Agent}
\date{\today}

\begin{document}
\maketitle

\begin{abstract}

\end{abstract}

\chapter{PENDAHULUAN}


\chapter{TINJAUAN PUSTAKA}


% Original content (commented out):
% \documentclass{article}
\usepackage[a4paper, top=2cm,right=1cm,left=1cm,bottom=0.5cm]{geometry}
\usepackage{url}

\usepackage{ifthen}
\usepackage{xcolor}

\usepackage{multicol}
\setlength{\columnsep}{1cm}
\setlength{\columnseprule}{1pt}
\usepackage[compact]{titlesec}

\usepackage{mathtools}
\usepackage{siunitx}
\sisetup{inter-unit-product=\ensuremath{{}\cdot{}}, exponent-product = \cdot}
\DeclareSIUnit\bar{bar}

\usepackage[version=4]{mhchem} 

\makeatletter
\newsavebox\myboxA
\newsavebox\myboxB
\newlength\mylenA
\newcommand*\overbar[2][0.75]{%
    \sbox{\myboxA}{$\m@th#2$}%
    \setbox\myboxB\null% Phantom box
    \ht\myboxB=\ht\myboxA%
    \dp\myboxB=\dp\myboxA%
    \wd\myboxB=#1\wd\myboxA% Scale phantom
    \sbox\myboxB{$\m@th\overline{\copy\myboxB}$}%  Overlined phantom
    \setlength\mylenA{\the\wd\myboxA}%   calc width diff
    \addtolength\mylenA{-\the\wd\myboxB}%
    \ifdim\wd\myboxB<\wd\myboxA%
       \rlap{\hskip 0.5\mylenA\usebox\myboxB}{\usebox\myboxA}%
    \else
        \hskip -0.5\mylenA\rlap{\usebox\myboxA}{\hskip 0.5\mylenA\usebox\myboxB}%
    \fi}
\makeatother
\renewcommand{\arraystretch}{1.2}

\makeatletter
\renewcommand*\env@matrix[1][\arraystretch]{%
  \edef\arraystretch{#1}%
  \hskip -\arraycolsep
  \let\@ifnextchar\new@ifnextchar
  \array{*\c@MaxMatrixCols c}}
\makeatother

\usepackage{esint}
\usepackage{setspace}
\usepackage{colortbl}
\usepackage{xcolor}
\usepackage{enumitem}

\usepackage{float}
\usepackage{fancyhdr}
\usepackage{lastpage,refcount, atbegshi}
\usepackage[hidelinks]{hyperref}
\AtBeginShipout{%
  \ifnum\value{page}=\number\numexpr\getpagerefnumber{LastPage}-2\relax
    \phantomsection\label{preLastPage}
  \fi}

\usepackage{polyglossia}

\setmainfont{David CLM}
\setsansfont{Miriam CLM}

\setdefaultlanguage{hebrew}
\setotherlanguage{english}  

\setmathrm{TeX Gyre Schola}

\makeatletter
\def\xpg@set@normalfont#1{%
  \letcs{\rmfamily}{#1@font@rm}%
  \letcs{\sffamily}{#1@font@sf}%
  \letcs{\ttfamily}{#1@font@tt}%
  \def\normalfont{\protect\xpg@select@fontfamily{#1}}%def instead of gdef
  \gdef\reset@font{\protect\normalfont}%
}
\addto\inlineextras@english{\xpg@set@normalfont{english}}
\addto\blockextras@english{\xpg@set@normalfont{english}}
\makeatother

\usepackage[style=english]{csquotes}

\usepackage[continuous]{pagenote}
\renewcommand*{\notesname}{הערות}
\renewcommand{\sectionname}{חלק}
\renewcommand{\notenuminnotes}[1]{{\textnormal#1. }}
\makepagenote

\renewcommand\text[1]{\textnormal{\textenglish{#1}}}

\parindent0pt
\pagestyle{fancy}
\fancyhf{}
\renewcommand{\headrulewidth}{0pt}
\setlength{\headsep}{0.3cm}
\fancyhead[L]{ עמוד {\thepage} מתוך \pageref{preLastPage}}
\fancyhead[R]{\ifthenelse{\value{page}=1}{\today}{}}

\def\imagewidth{0.9}

\newenvironment{cheatformula}[1][כותרת]{
    \begin{minipage}{\linewidth}
    \textbf{#1}:
}{
    \end{minipage}\\[2ex]
}

\newcommand{\cheatimage}[4][\imagewidth]{
    \begin{figure}[H]
        \centering
        \includegraphics[width=#1\linewidth]{#2}
        \caption{#3}
        \label{#4}
    \end{figure}
}

\newcommand*{\NameAndID}{%
    \par\noindent\makebox[2.5in]{\hrulefill} \hspace{0.5in}\makebox[2.0in]{\hrulefill}%
    \par\noindent\makebox[2.5in][r]{שם מלא}      \hspace{0.5in}\makebox[2.0in][r]{ת"ז}%
}%

\title{דף נוסחאות כללי}
\author{עידו פנג בנטוב}

\begin{document}

\makeatletter
\begin{center}
    {\NameAndID}\\[2ex]
    {\huge{\textbf{\@title}}}\\[2ex]
\end{center}
\makeatother

\begin{multicols*}{2}
\raggedcolumns

\section{משוואות מסדר שני בשני משתנים}
\begin{cheatformula}[מושגים]\\
    \pagenote{וואו איזה הערה}
    צורה כללית:
    $$au_{xx}+2bu_{xy}+cu_{yy}+du_{x}+eu_{y}+fu=0$$
    \textbf{הדיסקרימיננטה} $(\delta)$ של המשוואה מוגדרת כ:
    $$\delta\equiv b^{2}-ac$$
    ה\textbf{חלק העיקרי של המשוואה} מוגדר כ:
    $$au_{xx}+2bu_{xy}+cu_{yy}$$
\end{cheatformula}

\section{קריטריוני כניעה וכשל}
\begin{cheatformula}[קריטריון רנקין]
    $$\sigma_{\text{eq}}=\max_{}\left\{ \left|\sigma^{(i)}\right| \right\}$$
\begin{cheatformula}[קריטריון טרסקה]
    $$\sigma _{\text{eq}}=\max_{}\left\{ \left|\sigma^{(1)}-\sigma^{(2)}\right|,\, \left|\sigma^{(2)}-\sigma^{(3)}\right|,\, \left|\sigma^{(1)}-\sigma^{(3)}\right| \right\}$$
\end{cheatformula}
\end{cheatformula}
\begin{cheatformula}[קריטריון פון-מיזס]
    $$\sigma_{\text{eq}}=\dfrac{1}{\sqrt{ 2 }}\sqrt{ (\sigma^{(1)}-\sigma^{(2)})^{2}+(\sigma^{(1)}-\sigma^{(3)})^{2}+(\sigma^{(3)}-\sigma^{(2)})^{2} }$$
\end{cheatformula}
\begin{cheatformula}[מקדמי ביטחון $K$]
    $$\sigma_{\text{eq}}=\dfrac{\sigma _{y}}{K}$$
\end{cheatformula}
\begin{cheatformula}[מיכלי לחץ גליליים דקי דופן]
    $$\sigma_{rr}\approx 0,\, \quad \sigma_{zz}=\dfrac{1}{2}P \dfrac{R}{t},\,\quad \sigma_{\theta\theta}=P \dfrac{R}{t}$$
    אם המיכל נתון תחת פיתול $T$:
    $$\sigma_{\theta z}=\dfrac{Tr}{J}$$
\end{cheatformula}
\cheatimage{cyl.png}{החתכים שבעזרתם נמצאו המאמצים במיכל לחץ.}{cyl}

\section{כפיפה משופעת}
\begin{cheatformula}[אינרציה]
    \begin{gather*}
        X_{i} =\int _{A}x_{i} \, \mathrm{d}A\qquad
        I_{ij} =\int _{A} x_{i}x_{j} \, \mathrm{d}A
    \end{gather*}
\end{cheatformula}
\begin{cheatformula}[מרכז הכובד של התת-חתך]
       \pagenote{ישנו הבדל מאוד משמעותי בין $X_i$ ל-$Q_i$. בעוד $X_i$ מחושב על כל החתך, $Q_i$ מחושב על התת-חתך, כאשר $x_i$ נתוון ע"י מרכז המסה של כלל החתך $A$.}
       $$Q_{i}= \int _{\tilde{A}}x_{i} \, \mathrm{d}A=\bar{x}A$$
       כאשר $\bar{x}$ הוא מרכז הכובד של התת-חתך \textit{ביחס} למרכז הכובד של כלל החתך.
\end{cheatformula}
\begin{cheatformula}[מאמץ נורמלי בכפיפה משופעת]
    $$\sigma_{11}=\dfrac{N}{A}+\dfrac{({M}_{2} {I}_{22} +{M}_{3} {I}_{23} ){x}_{3} -({M}_{3} {I}_{33} +{M}_{2} {I}_{23} ){x}_{2} }{{I}_{22} {I}_{33} -({I}_{23})^{2} }$$
    במערכת ראשית:
    $$\sigma_{11}=\dfrac{N}{A}+\dfrac{{M}_{2}}{I_{33}}{x}_{3} -\dfrac{M_{3}}{I_{22}}{x}_{2} $$
    \textbf{ציר ניטראלי} יתקבל כאשר $\sigma_{11}=0$.
\end{cheatformula}
\begin{cheatformula}[מאמץ גזירה בכפיפה משופעת]
    \begin{align*}
        \tau & =-\dfrac{1}{t}\left( \dfrac{({V}_{3} {I}_{22} -V_{2}{I}_{23} ){Q}_{3} -({V}_{3} {I}_{23} -{V}_{2} {I}_{33} ){Q}_{2} }{{I}_{22} {I}_{33} -({I}_{23})^{2} } \right)
    \end{align*}
    במערכת ראשית:
    $$\tau=-\dfrac{1}{t}\left( \dfrac{{V}_{2} {Q}_{2} }{I_{22}}+\dfrac{{V}_{3} {Q}_{3} }{I_{33}} \right)$$
\end{cheatformula}
\begin{cheatformula}[משפט שטיינר]\\
    טנזור האינרציה $I_{ij}'$ לאחר ההעתקה של מערכת צירים $\Delta_{i},\Delta_{j}$ וטנזור אינרציה מקורי $I_{ij}$ נתון ע"י:
    \begin{align*}
        &I_{ij}'=I_{ij}+\Delta_{i}\Delta_{j}A-\Delta_{j}X_{i}-\Delta_{i}X_{j}\\
        &I_{ij}'=I_{ij}+\Delta _{i}\Delta _{j}A \qquad (X_i=X_j=0)
    \end{align*}
    ניתן לראות ממשפט זה שהרכיבים על האלכסון של טנזור האינרציה יהיו הכי קטנים כאשר הוא מחושב במרכז במסה.
\end{cheatformula}
\begin{cheatformula}[זרימת הגזירה]
    זרימת הגזירה הכוללת שנכנסת לצומת שווה לזרימת הגזירה הכוללת שיוצאת ממנה.
\end{cheatformula}

\section{שיטות אנרגיה}
\begin{cheatformula}[אנרגיית שינוי צורה]
    $$U=\dfrac{1}{2}\int _{V} \sigma_{ij}\varepsilon_{ij} \, \mathrm{d}V$$
    במקרה של קורה במתיחה:
    $$U=\int _{L} \dfrac{1}{2} \dfrac{N^{2}}{EA} \, \mathrm{d}{x}_{1}$$
    קורה בפיתול:
    $$U=\dfrac{1}{2}\int _{L} \dfrac{T^{2}}{GJ} \, \mathrm{d}{x}_{1} $$
    קורה בכפיפה:
    $$U =\dfrac{1}{2}\int _{L}  \dfrac{({M}_{2})^{2}}{E{I}_{33}} +\dfrac{({M}_{3})^{2}}{E{I}_{22}} \, \mathrm{d}{x}_{1}$$
\end{cheatformula}
\begin{cheatformula}[כוח והזזה מוכללים]\\
    \textbf{כוח מוכלל} (כוח מרוכז או מומנט) מסומן ב-$Q_i$.\\
   \textbf{ ההזזה המוכללת} (שקיעה/זווית) במקום ובכיוון בו פועל $Q_i$ מסומן ב-$q_i$.\\
   כוח והזזה מוכללים קשורים לינארית ע"י מטריצת ההיענות/הקשיחות:
   $$Q_{i}=K_{ij}q_{j} \quad \text{ or } \quad q_{i}=S_{ij}Q_{j}$$
   $S_{ij}$ משמעותו הוא הזזה של $q_{i}$ עקב כוח מוכלל $Q_{j}$ בגודל יחידה.
\end{cheatformula}
\begin{cheatformula}[משפט ההדדיות של בטי-מקסוול]
    $$S_{ij}=S_{ji}$$
\end{cheatformula}
\begin{cheatformula}[המשפט השני של קסטיליאנו]
    $$q_i=\dfrac{ \partial U }{ \partial Q_{i} }$$
\end{cheatformula}


\end{multicols*}
\newpage
\printnotes
\end{document}


\end{document}